% !TeX program = pdflatex
% =============================================================
% Public Economics -- Lecture 3
% Empirical Tools of Public Finance
% Based on: Gruber, J. "Public Finance and Public Policy", Ch. 3
% Austrian institutional and empirical context added for WU Vienna.
% =============================================================
\documentclass[11pt,aspectratio=169]{beamer}

\usetheme{metropolis}
\usepackage[utf8]{inputenc}
\usepackage[T1]{fontenc}
\usepackage{lmodern}
\usepackage{booktabs}
\usepackage{array}
\usepackage{textcomp}
\usepackage{siunitx}
\usepackage{tikz}
\usepackage{pgfplots}
\pgfplotsset{compat=1.18}
\usepackage{pgf-pie}
\usetikzlibrary{positioning,arrows.meta,calc}
\usepackage{fontawesome5}
% ----------------------------------------------------------------
% Colors: WU Vienna blue as primary accent
% ----------------------------------------------------------------
\definecolor{wublue}{HTML}{003D7C}
\definecolor{wulight}{HTML}{4F8FC4}
\definecolor{wugray}{HTML}{5C6670}
\definecolor{wured}{HTML}{B0202E}
\definecolor{wugreen}{HTML}{4C8C6B}
\definecolor{wuyellow}{HTML}{D9A441}
\setbeamercolor{palette primary}{bg=wublue,fg=white}
\setbeamercolor{frametitle}{bg=wublue,fg=white}
\setbeamercolor{progress bar}{fg=wublue,bg=wulight!30}
\setbeamercolor{alerted text}{fg=wured}

\setbeamertemplate{frame footer}{Empirical Tools of Public Finance}
\setbeamertemplate{footline}[frame number]

\title{Public Economics}
\subtitle{Lecture 3 \(\cdot\) Empirical Tools of Public Finance}
\author{Shushanik Margaryan}
\institute{WU Vienna University of Economics and Business}
\date{Winter Term}

% Source note command for footnote-style citations on data slides
\newcommand{\srcnote}[1]{{\footnotesize\color{wugray}Source: #1}}

\begin{document}

% =============================================================
\begin{frame}
  \titlepage
\end{frame}

 

% =============================================================
\begin{frame}{Outline}
  \tableofcontents
\end{frame}

% =============================================================
\section{3.1 Correlation Is Not Causation}

\begin{frame}{The Identification Problem}
  In the early nineteenth century, a cholera epidemic swept through Russia.
  \begin{itemize}
    \item The government sent doctors to the worst-affected areas.
    \item Peasants in one province observed a very high \textbf{correlation} between the number of doctors in an area and the incidence of cholera there.
    \item Relying on this fact, they banded together and murdered their doctors.
  \end{itemize}
  \vspace{0.5em}
  \begin{alertblock}{What went wrong}
    They correctly saw that doctors and illness moved together --- but the \emph{causal} arrow ran the other way: more illness brought more doctors, not the reverse.
  \end{alertblock}
  \srcnote{Gruber, \emph{Public Finance and Public Policy}, Ch.\ 3.1.}
\end{frame}

\begin{frame}{The Same Mistake Keeps Happening}
  \begin{itemize}
    \item September 2014: aid workers trying to help communities in Guinea fight an Ebola outbreak were attacked, and eight were killed. A local leader: \emph{``Wherever those people have passed, the communities have been hit by illness.''}
    \item 1988: a Harvard dean found that freshmen who had taken SAT-preparation courses scored, on average, \emph{63 points lower} than those who hadn't --- and concluded the courses were unhelpful.
    \item In fact: in 1998 prep courses were rare, so mostly the \emph{weakest} test-takers --- who needed the most help --- took them. The courses didn't cause worse scores; students likely to score worse were the ones taking the courses.
  \end{itemize}
  \vspace{0.3em}
  \centering
  \alert{Statistics calls this the \emph{identification problem}: given that two variables move together, how do you identify what's causing what?}\\
  \srcnote{Gruber, \emph{Public Finance and Public Policy}, Ch.\ 3.1.}
\end{frame}

\begin{frame}{Three (and Only Three) Explanations for Any Correlation}
  For any correlation between two variables $A$ and $B$:
  \begin{enumerate}
    \item \(A\) causes \(B\).
    \item \(B\) causes \(A\).
    \item Some third factor causes both.
  \end{enumerate}
  \vspace{0.5em}
  \begin{alertblock}{Why this matters for policy}
    Policy makers want to use empirical results to \emph{predict} how an intervention will change behavior. Correlation alone provides no predictive power --- prediction requires knowing the \emph{causal} link.
  \end{alertblock}
\end{frame}

% =============================================================
\section{3.2 Measuring Causation with Data We'd Like to Have: Randomized Trials}

\begin{frame}{The Gold Standard: Randomized Trials}
  \begin{block}{Randomized trial}
    A group of volunteers is randomly divided into a \textbf{treatment group}, which receives the intervention being studied, and a \textbf{control group}, which does not --- effectively, assignment by the flip of a coin.
  \end{block}
  \vspace{0.5em}
  \begin{itemize}
    \item Because assignment doesn't depend on anything about the subjects, the two groups are --- on average --- identical in \emph{every} respect except the treatment.
    \item Any remaining difference in outcomes can then be attributed to the treatment: this is what makes randomization the gold standard for estimating causal effects.
  \end{itemize}
\end{frame}

\begin{frame}{The Problem of Bias}
  \begin{block}{Bias}
    Any source of difference between treatment and control groups that is correlated with the treatment but is \emph{not due to} the treatment.
  \end{block}
  \vspace{0.3em}
  \begin{itemize}
    \item Cholera doctors, SAT prep, and breastfeeding studies all suffered from bias: treatment and control groups differed for reasons \emph{other than} the treatment itself.
    \item By definition, a properly randomized trial has \emph{no} such bias --- the only systematic difference between the groups is the treatment.
  \end{itemize}
  \vspace{0.3em}
  \centering
  \alert{This is why randomized trials are the benchmark against which every other empirical method is judged.}
\end{frame}

 

\begin{frame}{Randomized Trials in Public Finance: The Global Rise of RCTs}
 
  \begin{itemize}
    \item Randomized trials are far more common in medicine than in public finance --- but the last two decades have seen a surge, led by the \textbf{MIT Jameel Poverty Action Lab (J-PAL)}, founded 2003 by Abhijit Banerjee, Esther Duflo, and Sendhil Mullainathan.
    \item J-PAL's network has run 1,000+ randomized evaluations of social programs across 91 countries, in partnership with governments, NGOs, and firms.
  \end{itemize}
  \begin{alertblock}{Example: Indonesia's Raskin subsidized-rice program}
    Only a third of eligible households received their full rice entitlement, partly because households couldn't prove their eligibility. Researchers randomly gave 378 villages an ID card stating a household's entitlement, versus 194 control villages without one. Card villages saw a sizeable rise in rice actually received --- without more going to ineligible households. Indonesia scaled the card nationwide in 2013.
  \end{alertblock}
  \srcnote{Banerjee, Hanna \& Olken, cited in Gruber, \emph{Public Finance and Public Policy}, Ch.\ 3.2.}
\end{frame}

\begin{frame}{A European Randomized Trial: Job-Search Counseling in France}
  
  \begin{itemize}
    \item Public-finance RCTs aren't only a developing-economy story. Behaghel, Crépon \& Gurgand (2014) ran a large controlled experiment across roughly 235 French local labor markets (\(\approx\)44{,}000 job seekers).
    \item Job seekers were randomly assigned to intensive counseling from the \textbf{public} employment service, from \textbf{private} providers, or to a control group receiving standard services.
  \end{itemize}
  \begin{alertblock}{Finding}
    Intensive counseling raised job-exit rates by roughly \textbf{15--35\%} relative to control --- and public-service counseling was about twice as effective as private-provider counseling in the first six months.
  \end{alertblock}

  \srcnote{Behaghel, Crépon \& Gurgand, \emph{AEJ: Applied Economics} 6(4), 2014; Crépon et al., \emph{QJE}, 2013.}
\end{frame}

\begin{frame}{Why We Need to Go Beyond Randomized Trials}
  \begin{itemize}
    \item \textbf{Cost and time}: large trials are expensive and slow to plan and run.
    \item \textbf{Ethics}: 
    \item \textbf{External validity}:
    \item \textbf{Attrition}: if people leave the trial non-randomly (e.g., those with negative side effects drop out), the remaining sample is no longer representative.
  \end{itemize}

\end{frame}

% =============================================================
\section{3.3 Estimating Causation with the Data We Actually Get}

\begin{frame}{Observational Data}
  \begin{block}{Observational data}
    Data generated by individual behavior observed in the real world, not in the context of a deliberately designed experiment.
  \end{block}
  \vspace{0.5em}
  Most of the time, this is all we have. The rest of this lecture covers three approaches researchers use to get \emph{as close as possible} to the gold standard using observational data:
  \begin{enumerate}
    \item Time series analysis
    \item Cross-sectional regression analysis
    \item Quasi-experiments  
  \end{enumerate}
\end{frame}

\begin{frame}{Time Series Analysis}
  \begin{block}{Time series analysis}
    Documenting the correlation between two variables of interest \emph{over time}.
  \end{block}
  \vspace{0.5em}
  \begin{itemize}
    \item Simplest possible empirical approach: plot both series across years and see whether they move together.
    \item \textbf{Problem}: when both series follow a slow-moving trend, almost \emph{anything} else that trends over the same period will look correlated with it --- the correlation may reflect a shared trend, not causation.
  \end{itemize}
\end{frame}

\begin{frame}{Spurious correlations}
	% TODO: \usepackage{graphicx} required
	\begin{figure}
		\centering
		\includegraphics[width=0.8\linewidth]{spurious}

	\end{figure}
	
\end{frame}
 
\begin{frame}{When Is Time Series Analysis Useful?}
 
  
  \begin{alertblock}{A sharp break: Ireland's 2002 plastic-bag levy}
   \end{alertblock}
   
   \begin{itemize}
   	\item  In March 2002, Ireland introduced a \texteuro{}0.15 levy per plastic shopping bag --- one precisely dated policy change. 
   	\item Convery, McDonnell \& Ferreira (2007) document a reduction in plastic-bag use ``in the order of 90\%'' at retail outlets; a widely-cited companion figure (Irish government) puts per-capita bag consumption at 328/year before the levy versus 21/year after (\(\approx\)95\%).
   
   	\item The drop happens \emph{exactly} at the levy's start date, not gradually --- hard to explain with a coincidental third factor. That is what makes a sharp break more convincing than a slow trend.\\
 
   \end{itemize} 	
     	\srcnote{Convery, McDonnell \& Ferreira, ``The Most Popular Tax in Europe? Lessons from the Irish Plastic Bags Levy,'' \emph{Environmental and Resource Economics} 38(1), 2007
}
\end{frame}

\begin{frame}{Cross-Sectional Regression Analysis}
  \begin{block}{Cross-sectional regression analysis}
    A statistical method for assessing the relationship between two (or more) variables \emph{while holding other factors constant}, comparing many individuals at \emph{one point in time}.
  \end{block}
  \vspace{0.3em}
  \begin{itemize}
    \item In its simplest form (\emph{bivariate regression}), it formalizes correlation: it finds the line that \emph{best fits} a scatter of data points, and reports its slope.
    \item The \textbf{dependent variable} is what you want to explain; the \textbf{independent variable(s)} are what you think explain it.
  \end{itemize}
\end{frame}

\begin{frame}{Regression fit}
\begin{figure}
	\centering
	\includegraphics[width=0.8\linewidth]{ols}
 
\end{figure}
\faHandPointer \href{https://setosa.io/ev/ordinary-least-squares-regression/}{Interactive OLS visualization (setosa.io)}

\end{frame}
% TODO: \usepackage{graphicx} required

 



\begin{frame}{Problems with Cross-Sectional Regression: Bias Again}
 \begin{itemize}
 	\item   With real data we don't observe ``preference for leisure'' directly --- but the same bias risk remains: people with a low taste for work may simply be both lower earners \emph{and} higher benefit recipients, with no causal link running from benefits to hours at all.

 \end{itemize}

\end{frame}

\begin{frame}{Control Variables}
  \begin{block}{Control variables}
    Variables included in a regression to account for differences between comparison groups that could otherwise bias the estimate.
  \end{block}
  \vspace{0.3em}
  \begin{itemize}
    \item If we had a perfect measure of ``preference for leisure,'' we could compare only people with the \emph{same} preference --- removing that source of bias entirely.
  
    \item The variables we really want (intrinsic preferences, motivation, ability) are usually impossible to measure. 
    \item We approximate them with proxies --- age, education, past employment --- which only partially capture the underlying difference. 
    \item Control variables reduce bias; they rarely eliminate it.
  \end{itemize}
 
\end{frame}

\begin{frame}{Quasi-Experiments}
  \begin{block}{Quasi-experiment}
    A situation that arises naturally when a change in the economic environment (typically, a policy change) creates treatment and control groups that are \emph{nearly identical} except for that change --- letting an outside event do the randomizing for us.
  \end{block}
  \vspace{0.3em}
 
  Austria's federal structure --- nine \emph{Länder}, each able to set some of its own policy parameters --- is a natural source of quasi-experiments: a reform in one Land, with neighbors unaffected, creates a ready-made treatment/control comparison.
  

\end{frame}

\begin{frame}
	There are many empirical strategies and designs that utilize quasi-experiments to estimate the causal effect of a policy. Among the most commonly used ones are: 
	
	\begin{itemize}
		\item Difference-in-differences
		\item Instrumental variable approaches 
		\item Regression discontinuity 
	\end{itemize}
 
 
	\faHandPointer  \href{https://nickchk.com/causalgraphs.html}{\textit{Causal Inference Animated Plots}}
\end{frame}
\begin{frame}{Illustrative Example (Numbers Made Up)}
  \footnotesize
  Suppose two neighboring Bundesländer, \emph{Vorarlberg} and \emph{Tirol}, both have single-parent Mindestsicherung recipients. In 2023, Vorarlberg cuts its monthly benefit; Tirol's stays unchanged.
  \begin{center}
  \begin{tabular}{@{}lccc@{}}
    \toprule
     & \textbf{2022} & \textbf{2024} & \textbf{Difference} \\
    \midrule
    \textbf{Vorarlberg} (benefit) & \texteuro{}1{,}000 & \texteuro{}800 & $-$\texteuro{}200 \\
    Vorarlberg (hours/year) & 1{,}000 & 1{,}200 & $+$200 \\
    \midrule
    \textbf{Tirol} (benefit) & \texteuro{}1{,}000 & \texteuro{}1{,}000 & \texteuro{}0 \\
    Tirol (hours/year) & 1{,}050 & 1{,}100 & $+$50 \\
    \bottomrule
  \end{tabular}
  \end{center}
  \vspace{0.3em}
  Vorarlberg's raw before/after change (+200 hours) mixes the \emph{benefit cut} with anything else moving nationally (e.g., an economic upswing) --- exactly the time series problem. Tirol didn't experience the cut, but \emph{did} experience the same national trend.
\end{frame}

 

\begin{frame}{Problems with Quasi-Experimental Analysis}
  \begin{itemize}
    \item Diff-in-diff assumes the \textbf{parallel-trends} assumption: absent the reform, treatment and control would have moved \emph{the same way}.
    \item If the national trend hit Vorarlberg and Tirol \emph{differently} (e.g., their labor markets are exposed to different industries), the ``bias'' terms are not equal --- and simple subtraction no longer isolates the true treatment effect.
    \item Quasi-experiments can never be \emph{certain} they've purged all bias, the way a true randomized trial can. Researchers argue their case both intuitively (the groups look comparable) and statistically (checking the result against alternative control groups).
  \end{itemize}

\end{frame}

% =============================================================
\section{Quiz: Correlation or Causation?}

\begin{frame}{How to Play}
  \centering
  \Large \faQuestionCircle
  \vspace{0.5em}
  \normalsize
  \begin{block}{The rules}
    Each round: a scenario, then a multiple-choice question. Pick your answer --- then click to reveal.
  \end{block}
  \begin{alertblock}{The question to always ask}
    Is there anything \emph{other than} the variable of interest that differs, systematically, between the groups being compared --- and that could \emph{also} move the outcome?
  \end{alertblock}
\end{frame}

\begin{frame}{Round 1: Police and Crime}
  \small
  A city with more police officers per capita also has a \emph{higher} recorded crime rate than cities with fewer police. A researcher regresses crime rates on police-per-capita and reports \(\hat\beta > 0\).
  \vspace{0.5em}
  \textbf{What best explains \(\hat\beta > 0\)?}
  \begin{enumerate}
    \item[\textbf{A}] Police presence causes crime
    \item[\textbf{B}] Reverse causality: high crime causes more police deployment
    \item[\textbf{C}] Pure coincidence
  \end{enumerate}
  \pause
  \vspace{0.3em}
  \begin{alertblock}{\faCheckCircle\ \ Answer: B}
    Cities allocate \emph{more} police to areas that \emph{already} have more crime. Same logic as the cholera doctors from \S3.1.
  \end{alertblock}
\end{frame}

\begin{frame}{Round 2: Ice Cream and Drowning}
  \small
  Across a summer season, towns with higher weekly ice-cream sales also record more drowning deaths. \(\hat\beta > 0\) from a regression of drownings on ice-cream sales.
  \vspace{0.5em}
  \textbf{What best explains \(\hat\beta > 0\)?}
  \begin{enumerate}
    \item[\textbf{A}] Ice cream causes risky swimming behavior
    \item[\textbf{B}] Drowning risk causes people to crave ice cream
    \item[\textbf{C}] A third factor --- hot weather --- drives both up independently
  \end{enumerate}
  \pause
  \vspace{0.3em}
  \begin{alertblock}{\faCheckCircle\ \ Answer: C}
    Hot weather raises both ice-cream sales \emph{and} swimming (hence drownings). Neither variable causes the other --- the classic textbook confounder.
  \end{alertblock}
\end{frame}

\begin{frame}{Round 3: Schooling and Wages}
  \small
  A researcher regresses hourly wages on years of schooling across a large sample of workers and finds a positive \(\hat\beta\), with no control variables included.
  \vspace{0.5em}
  \textbf{Is \(\hat\beta\) unbiased?}
  \begin{enumerate}
    \item[\textbf{A}] Yes --- more schooling clearly pays off
    \item[\textbf{B}] No --- likely omitted-variable bias from unobserved ability/motivation
    \item[\textbf{C}] No --- reverse causality (wages cause schooling)
  \end{enumerate}
  \pause
  \vspace{0.3em}
  \begin{alertblock}{\faCheckCircle\ \ Answer: B}
    People with more underlying ability or drive tend to get \emph{both} more schooling \emph{and} higher wages --- even if schooling itself had zero causal effect. \(\hat\beta\) is likely biased upward.
  \end{alertblock}
\end{frame}

\begin{frame}{Round 4: Adding a Control}
  \small
  The same researcher adds \emph{parental income} as a control variable. The coefficient on schooling shrinks but stays positive and significant.
  \vspace{0.5em}
  \textbf{Is \(\hat\beta\) unbiased now?}
  \begin{enumerate}
    \item[\textbf{A}] Yes --- the bias is now fully removed
    \item[\textbf{B}] Not necessarily --- other omitted factors (ability, work ethic, school quality) remain uncontrolled
    \item[\textbf{C}] No --- adding controls always makes bias worse
  \end{enumerate}
  \pause
  \vspace{0.3em}
  \begin{alertblock}{\faCheckCircle\ \ Answer: B}
    Parental income proxies for only \emph{part} of what drives both schooling and wages. Control variables reduce bias only for what they actually capture.
  \end{alertblock}
\end{frame}

\begin{frame}{Round 5: A Benefit Cut, Two Regions}
  \small
  Region A cuts a welfare benefit; neighboring Region B leaves it unchanged. A researcher compares the before/after change in hours worked in A to the before/after change in B (a difference-in-differences design) and reports a large positive effect on labor supply.
  \vspace{0.5em}
  \textbf{Is \(\hat\beta\) unbiased?}
  \begin{enumerate}
    \item[\textbf{A}] Yes --- diff-in-differences always removes bias
    \item[\textbf{B}] Only if Region A and B were on the \emph{same} underlying trend absent the cut
    \item[\textbf{C}] No --- this design can never be trusted
  \end{enumerate}
  \pause
  \vspace{0.3em}
  \begin{alertblock}{\faCheckCircle\ \ Answer: B}
    This is the \textbf{parallel-trends} assumption. If Region A's economy was already diverging from Region B's for unrelated reasons, simple subtraction no longer isolates the causal effect.
  \end{alertblock}
\end{frame}

\begin{frame}{Round 6: Deworming in Kenya}
  \small
  Miguel \& Kremer (2004) randomized a school-based deworming program across dozens of primary schools in western Kenya, phasing treatment in school-by-school rather than child-by-child.
  \vspace{0.5em}
  \textbf{Why randomize by \emph{school}, not by individual child, within a school?}
  \begin{enumerate}
    \item[\textbf{A}] It's cheaper to administer
    \item[\textbf{B}] To capture \emph{spillovers}: untreated children in a treated school also benefit from reduced worm transmission, which individual-level randomization would hide
    \item[\textbf{C}] It doesn't matter --- either design gives the same answer
  \end{enumerate}
  \pause
  \vspace{0.3em}
  \begin{alertblock}{\faCheckCircle\ \ Answer: B}
    Randomizing by individual \emph{within} a school would understate the true effect --- untreated ``controls'' next to treated children aren't a clean comparison when the treatment has spillovers.
  \end{alertblock}
  \srcnote{Miguel \& Kremer, ``Worms: Identifying Impacts on Education and Health in the Presence of Treatment Externalities,'' \emph{Econometrica} 72(1), 2004.}
\end{frame}

% =============================================================
\section{Structural Modeling}

\begin{frame}{Reduced-Form vs.\ Structural Estimates}
	\begin{itemize}
		\item Structural estimates: Builds a theoretical model of individual behavior and then
		estimates the parameters of the model. Estimates of the features that drive individual
		decisions, such as income and substitution effects or parameters of the utility
		function.
		\item Reduced form estimates: Measures of the total impact of an independent variable on
		a dependent variable, without de- composing the source of that behavior response in
		terms of underlying parameters of the utility functions
		
		\item Reduced form estimates are more transparent and convincing but structural
		estimates are more directly useful to make predictions for alternative policies
	\end{itemize}

\end{frame}

% =============================================================
 

\begin{frame}[standout]
  Thank you!\\
  \vspace{0.5em}
  \small Questions and discussion
\end{frame}

\end{document}
